\section{M6 --- Correlação com Eventos e Comparação Setorial}
\label{sec:m6}

\subsection{Validação narrativa}
Complementando a avaliação quantitativa (Seção~\ref{sec:m5}), as anomalias detectadas no
teste (limiar estático p95) são cruzadas com uma linha do tempo curada de eventos
econômicos e políticos brasileiros de 2010--2024 (\texttt{src/events.py}; ADR-0008). Cada
evento é classificado como \emph{sistêmico} (afeta todos os ativos) ou \emph{específico} de
um \emph{ticker}. Considera-se um evento ``coberto'' quando há ao menos uma anomalia dentro
de uma tolerância de $\pm$\texttt{window\_size} ($\pm$30 dias), justificada pelo fato de a
unidade de detecção ser a janela de 30 passos.

A cobertura narrativa global é de $\approx$50\%, mas o número agregado esconde o resultado
relevante --- o \textbf{casamento entre evento e ativo afetado}:
\begin{itemize}[noitemsep]
  \item \textbf{AMER3} cobre 9 de 10 eventos a ela associados, incluindo a fraude contábil
        (jan/2023) e a recuperação judicial --- o sinal mais forte do \emph{dataset}.
  \item PETR4, VALE3 e ITUB4 cobrem cerca de 2 eventos cada (de 6--8): seus modelos disparam
        sobretudo em eventos \emph{sistêmicos} e nos choques setoriais de maior magnitude,
        permanecendo silenciosos em eventos de impacto difuso --- coerente com um detector
        calibrado para a cauda do erro (p95).
\end{itemize}

\begin{figure}[h]
\centering
\includegraphics[width=\textwidth]{m6_amer3_detalhe}
\caption{AMER3 (2022--2024): erro de reconstrução, janelas anômalas (pontos vermelhos) e
eventos (linhas pontilhadas). O pico de erro de jan/2023 coincide com a fraude contábil e o
pedido de recuperação judicial.}
\label{fig:m6-amer3}
\end{figure}

\subsection{Comparação setorial}
Com um modelo por ativo, comparam-se as taxas de marcação entre setores
(Tabela~\ref{tab:setorial}). As taxas reproduzem o achado da M4: PETR4 (energia) e VALE3
(mineração) são estáveis ($\approx$3\%), enquanto AMER3 (varejo) exibe regime de
volatilidade prolongada (26{,}7\% no estático, recalibrado para 14{,}3\% no dinâmico) e
ITUB4 (financeiro) fica em posição intermediária. A diferença reflete o \emph{ativo}, não
necessariamente o \emph{setor} ($n=1$ por setor): trata-se de estudo de caso qualitativo,
não de inferência estatística (ADR-0006).

\begin{table}[h]
\centering
\caption{Taxa de janelas anômalas no teste (2020--2024) por ativo/setor.}
\label{tab:setorial}
\begin{tabular}{llccc}
\toprule
\textbf{Ticker} & \textbf{Setor} & \textbf{Limiar p95} & \textbf{Estático (\%)} & \textbf{Dinâmico (\%)} \\
\midrule
PETR4.SA & Energia/commodities & 0{,}124 & 2{,}9  & 2{,}9  \\
VALE3.SA & Mineração           & 0{,}082 & 2{,}7  & 3{,}4  \\
AMER3.SA & Varejo              & 0{,}077 & 26{,}7 & 14{,}3 \\
ITUB4.SA & Financeiro          & 0{,}088 & 11{,}0 & 4{,}9  \\
\bottomrule
\end{tabular}
\end{table}

\subsection{Contágio sistêmico: o teste da COVID}
Um evento sistêmico deve aparecer em \emph{todos} os ativos. No \emph{crash} da COVID
(fev--mai/2020), os quatro modelos sinalizam anomalia de forma concentrada no período
(PETR4: 35; VALE3: 33; ITUB4: 54; AMER3: 56 janelas), confirmando o contágio entre setores.
Em contraste, eventos setor-específicos --- Brumadinho ($\to$ VALE3), fraude Americanas
($\to$ AMER3), intervenção na Petrobras ($\to$ PETR4) --- isolam-se nos ativos afetados.
A combinação dos dois comportamentos evidencia que os modelos, treinados por ativo, captam
tanto choques de mercado amplos quanto eventos idiossincráticos --- o objetivo de
generalização cruzada declarado na proposta.

\begin{figure}[h]
\centering
\includegraphics[width=\textwidth]{m6_covid_contagio}
\caption{\emph{Crash} da COVID (fev--mai/2020): os quatro ativos, de setores distintos,
sinalizam anomalias concentradas no mesmo período, confirmando o contágio sistêmico.}
\label{fig:m6-covid}
\end{figure}

\paragraph{Limitações em aberto.} O estudo de sensibilidade da janela dinâmica
(\texttt{dynamic\_window} $\in \{126, 252, 504\}$) previsto no ADR-0005 e a calibração das
magnitudes de injeção relativas a $\sigma$ (ADR-0006) permanecem pendentes e são os próximos
refinamentos antes do fechamento do artigo.
